\chapter{Extracting the Pauli observables}
\label{chap:qpbt-extraction}
% Paper subsection label sec:separating kept for cross-reference fidelity.
\label{sec:separating}
% Source: references/qpbt-paper/14_analysis_of_the_pauli_basis_test.tex,
% lines 1415--1877 (subsection ``Pulling the X and Z measurements apart''
% and the closing proof of the soundness theorem).

This chapter completes the analysis of the Pauli basis test~\cite{Ji2020MIPStar}. Throughout, $(q,m,d)$ is an admissible parameter tuple (Definition~\ref{def:admissible}), $q = 2^t$, and $M = 2^m$ denotes the number of $q$-ary qudits tested for per player. (The letter $M$ also denotes measurement operators, always carrying super- and subscripts; the meaning is determined by context, as in the source.) We fix a projective strategy for the Pauli basis test game $\mathfrak{G}^{\mathrm{Pauli}}_{(q,m,d)}$ that succeeds with probability at least $1-\eps$, in the setting of Chapter~\ref{chap:qpbt-observables}: the shared state $\ket{\psi}$ on registers $\mathrm{A}$, $\mathrm{B}$ is assumed padded with sufficiently many ancilla qubits in the state $\ket{0}$ (Lemma~\ref{lem:projective-strategy-setup} and the convention following it), and $\ket{\hat\psi}$ denotes the expanded state on the six registers $\mathrm{A}\mathrm{A}'\mathrm{A}''\mathrm{B}\mathrm{B}'\mathrm{B}''$, obtained by adjoining to each player the two halves of $M$ pairs $\ket{\mathrm{EPR}_q}$ (Definition~\ref{def:expanded-state}). Unless stated otherwise, every $\approx_\delta$ and $\simeq_\delta$ statement below is taken with respect to the state $\ket{\hat\psi}$, over the question distribution named in the statement, with register placements indicated by subscripts; ``all symmetric equivalents'' refers to the register interchanges of Definition~\ref{def:symmetric-equivalents}.

We use the following objects from the preceding chapters: the combined point projectors $\hat{M}^{(\mathrm{Point},W),u}_a$ acting on the expanded space (Definition~\ref{def:expanded-point-measurement}); the projective polynomial-pair measurement $\{\hat{S}_{g_X,g_Z}\}$ of Lemma~\ref{lem:qld-4-7}, with outcomes pairs of polynomials in $\mathrm{ideg}_{d,m}(\F_q)$ and error function
\[
	\delta_S(\eps,m,d,q) = a(md)^a\bigl(\eps^b + q^{-b} + 2^{-bmd}\bigr)
\]
for universal constants $a > 1$, $0 < b < 1$; the decoding map $\mathrm{Dec} = \mathrm{Dec}_{\F_q}\colon \mathrm{ideg}_{d,m}(\F_q) \to \F_q^{M}$ of the low-degree code (Definition~\ref{def:decoding-map}, applied with $H = \F_q$; see Remark~\ref{rem:qld-decoding-identity}), which maps a polynomial $g$ to the string of its evaluations on the hypercube $\{0,1\}^m$; and the low-degree encoding $h \mapsto g_h$ together with the vectors $\mathrm{ind}_m(u) \in \F_q^M$, characterized by $h \cdot \mathrm{ind}_m(u) = g_h(u)$ for all $u \in \F_q^m$ \eqref{eq:low-degree-encoding-definition}. Both maps are $\F_q$-linear, and $\mathrm{Dec}(g_h) = h$ for every $h \in \F_q^{M}$. We call a polynomial $g \in \mathrm{ideg}_{d,m}(\F_q)$ an \emph{encoding} if $g = g_h$ for some $h \in \F_q^{M}$ or, equivalently, if $g$ is multilinear; since a multilinear polynomial is determined by its values on $\{0,1\}^m$, the encodings are exactly the $g$ with $g_{\mathrm{Dec}(g)} = g$ as polynomial representatives, and
\[
	\mathrm{Dec}(g) \cdot \mathrm{ind}_m(u) = g(u)
	\qquad \text{for every encoding $g$ and every } u \in \F_q^m\;.
\]
For a $g$ that is not an encoding the two sides of this identity are the evaluations of the distinct polynomial representatives $g_{\mathrm{Dec}(g)}$ and $g$; when $d \leq q-1$ they consequently disagree at some $u$, but when $d \geq q$ the two representatives may induce the same function on $\F_q^m$, so pointwise validity of the identity does not characterize the encodings; see Remark~\ref{rem:qld-decoding-identity}.

% Local correction of the source (lines 1419 and 1492 of
% 14_analysis_of_the_pauli_basis_test.tex): the source invokes its boolean
% decoding map but describes it as returning all hypercube evaluations, and
% asserts the decoder--encoder identity for every g of individual degree at
% most d. The remark below records the corrected statements.
\begin{remark}[Choice of decoding map and the decoder--encoder identity]\label{rem:qld-decoding-identity}
	The source states the constructions of this chapter in terms of its boolean decoding map (the case $H = \{0,1\}$ of Definition~\ref{def:decoding-map}) while describing that map as returning the string of \emph{all} evaluations of $g$ on the hypercube. The proofs require the latter reading, $H = \F_q$, which is adopted throughout this chapter: the linearity of $\mathrm{Dec}$ and the identity $\mathrm{Dec}(g_h) = h$ for all $h \in \F_q^{M}$, both of which fail for the boolean map, are used in the symmetry step of the proof of Lemma~\ref{lem:qld-construct-the-paulis} and in the conjugation identity \eqref{eq:qld-unitary-6}. The source further asserts the identity $\mathrm{Dec}(g) \cdot \mathrm{ind}_m(u) = g(u)$ for every $g \in \mathrm{ideg}_{d,m}(\F_q)$. Since the left-hand side equals $g_{\mathrm{Dec}(g)}(u)$, the identity, taken for all $u \in \F_q^m$, states that the multilinear representative $g_{\mathrm{Dec}(g)}$ and $g$ induce the same function on $\F_q^m$. Every encoding satisfies it, because then $g_{\mathrm{Dec}(g)} = g$ as polynomial representatives; a polynomial that is not an encoding need not: for $m = 1$, $d \geq 2$, and $q > 2$, the polynomial $g(x) = x^2$ has $\mathrm{Dec}(g) \cdot \mathrm{ind}_1(u) = u$, which differs from $g(u) = u^2$ at every $u \notin \{0,1\}$. When $d \leq q-1$ the failure at some $u$ characterizes the polynomials that are not encodings, since evaluation is injective on representatives of individual degree at most $q-1$ (Definition~\ref{def:polynomials-degree}); when $d \geq q$ it does not: for the admissible tuple $(q,m,d) = (8,1,8)$, the polynomial $g(x) = x^8$ is not an encoding, yet $\mathrm{Dec}(g) \cdot \mathrm{ind}_1(u) = u = u^8 = g(u)$ for every $u \in \F_8$. Accordingly, the proofs below invoke the identity only for encodings; in the proof of Lemma~\ref{lem:qld-construct-the-paulis} the contribution of the remaining outcomes is controlled by the bound \eqref{eq:qld-nonencoding-mass} on the mass that the measurement of Lemma~\ref{lem:qld-4-7} places on polynomials that are not encodings, while in the conjugation identity \eqref{eq:qld-unitary-6} the quantity $\mathrm{Dec}(g) \cdot \mathrm{ind}_m(u)$ cancels exactly.
\end{remark}

\section{The pulled-apart measurements and observables}

The measurement $\{\hat{S}_{g_X,g_Z}\}$ reads out an $X$-polynomial and a $Z$-polynomial simultaneously. We now separate the two readings again, this time as commuting actions on distinct registers: the polynomial outcome is measured on registers $\mathrm{A}\mathrm{A}'$ (resp.\ $\mathrm{B}\mathrm{B}'$), while a genuine Pauli measurement acts on the EPR register $\mathrm{A}''$ (resp.\ $\mathrm{B}''$).

\begin{definition}[Single-basis marginals of the polynomial-pair measurement]\label{def:s-w-marginals}
	\uses{lem:qld-4-7}
	Let $\{\hat{S}_{g_X,g_Z}\}$ be the projective measurement of Lemma~\ref{lem:qld-4-7}. For every $g \in \mathrm{ideg}_{d,m}(\F_q)$ define
	\[
		\hat{S}^X_g = \sum_{g_Z} \hat{S}_{g,g_Z}
		\qquad \text{and} \qquad
		\hat{S}^Z_g = \sum_{g_X} \hat{S}_{g_X,g}\;,
	\]
	the sums ranging over $\mathrm{ideg}_{d,m}(\F_q)$. Since $\{\hat{S}_{g_X,g_Z}\}$ is projective, its elements are mutually orthogonal projectors; hence each $\hat{S}^W_g$, for $W \in \{X,Z\}$, is a projector, and $\{\hat{S}^W_g\}_{g}$ is a projective measurement with outcomes in $\mathrm{ideg}_{d,m}(\F_q)$.
\end{definition}

\begin{definition}[Pauli dot-product projectors]\label{def:tau-dot-product-projector}
	\uses{def:generalized-pauli}
	For $W \in \{X,Z\}$ and $h \in \F_q^{M}$ let $\tau^W_h$ denote the projector corresponding to measuring $(\C^q)^{\ot M}$ in the $W$ basis and obtaining the outcome $h$ (Definition~\ref{def:generalized-pauli}). For all $W \in \{X,Z\}$, $\tilde{u} \in \F_q^{M}$, and $a \in \F_q$, define the projector
	\begin{equation}\label{eq:def-tauwu}
		\tau^{W}_{a}(\tilde{u}) = \sum_{h \in \F_q^{M} \,:\, h \cdot \tilde{u} = a} \tau^W_{h}\;.
	\end{equation}
	For each fixed $\tilde{u}$, the family $\{\tau^W_a(\tilde{u})\}_{a \in \F_q}$ is a projective measurement on $(\C^q)^{\ot M}$.
\end{definition}

\begin{definition}[Pulled-apart Pauli basis measurements]\label{def:tilde-m-measurement}
	\uses{def:s-w-marginals, def:tau-dot-product-projector, def:decoding-map}
	For all $W \in \{X,Z\}$ and $\tilde{u} \in \F_q^{M}$ define the measurement $\{\tilde{M}^{W,\tilde{u}}_a\}_{a \in \F_q}$ by
	\begin{equation}\label{eq:tilde_M}
		\tilde{M}^{W,\tilde{u}}_a = \sum_{g} \hat{S}^W_g \ot \tau_{(\mathrm{Dec}(g) \cdot \tilde{u}) - a}^W(\tilde{u})\;,
	\end{equation}
	the sum ranging over $g \in \mathrm{ideg}_{d,m}(\F_q)$. If $\hat{S}^W_g$ is viewed as an operator acting on registers $\mathrm{A}\mathrm{A}'$ (resp.\ $\mathrm{B}\mathrm{B}'$), then $\tilde{M}^{W,\tilde{u}}_a$ acts on registers $\mathrm{A}\mathrm{A}'\mathrm{A}''$ (resp.\ $\mathrm{B}\mathrm{B}'\mathrm{B}''$), the $\tau^W$ factor acting on $\mathrm{A}''$ (resp.\ $\mathrm{B}''$). Each family $\{\tilde{M}^{W,\tilde{u}}_a\}_{a \in \F_q}$ is a projective measurement, because the $\hat{S}^W_g$ and the $\tau^W$ projectors are projective and, for each $g$, the projectors $\tau^W_{(\mathrm{Dec}(g)\cdot\tilde{u})-a}(\tilde{u})$ sum to the identity as $a$ ranges over $\F_q$.
\end{definition}

\begin{definition}[Pulled-apart Pauli observables]\label{def:tilde-w-observables}
	\uses{def:tilde-m-measurement, def:dual-self-dual-normal-basis, def:binary-representation, def:subfield-trace}
	Let $\{e_j\}_{j \in \{1,\ldots,t\}}$ be the self-dual normal basis for $\F_q$ over $\F_2$ (Definition~\ref{def:dual-self-dual-normal-basis}) that is fixed once and for all in Definition~\ref{def:binary-representation}, and let $\tr = \tr_{q \to 2} \colon \F_q \to \F_2$ be the subfield trace. For all $W \in \{X,Z\}$, $\tilde{u} \in \F_q^M$, and $j \in \{1,\ldots,t\}$ define the operator
	\begin{equation}\label{eq:def-tildewj}
		\tilde{W}^j(\tilde{u}) = \sum_{a \in \F_q} (-1)^{\tr(e_j a)} \, \tilde{M}_a^{W,\tilde{u}}\;,
	\end{equation}
	acting on the same registers as $\tilde{M}^{W,\tilde{u}}_a$.
\end{definition}

% In the source the following properties are stated inline with the
% definition of the observables; they are recorded here as a lemma.
\begin{lemma}[Product form and Pauli relations of the pulled-apart observables]\label{lem:tildew-product-form}
	\uses{def:tilde-w-observables, def:s-w-marginals, def:generalized-pauli, def:decoding-map, lem:qld-4-7}
	Let $\tau^W(\cdot)$ denote the generalized Pauli $W$ observable (Definition~\ref{def:generalized-pauli}). For all $W \in \{X,Z\}$, $\tilde{u} \in \F_q^M$, and $j \in \{1,\ldots,t\}$,
	\begin{equation}\label{eq:tildewj-product-form}
		\tilde{W}^j(\tilde{u})
		= \Big( \sum_{g_X,g_Z} (-1)^{\tr(e_j (\mathrm{Dec}(g_W) \cdot \tilde{u}))} \, \hat{S}_{g_X,g_Z} \Big) \ot \tau^W(e_j \tilde{u})\;,
	\end{equation}
	where $g_W$ denotes $g_X$ if $W = X$ and $g_Z$ if $W = Z$. In particular each $\tilde{W}^j(\tilde{u})$ is Hermitian and squares to the identity. Moreover, the pulled-apart observables satisfy exactly the twisted commutation relations of the generalized Pauli observables $\tau^X(e_j \tilde{u})$ and $\tau^Z(e_{j'} \tilde{v})$: for all $j,j' \in \{1,\ldots,t\}$ and $\tilde{u},\tilde{v} \in \F_q^M$,
	\begin{equation}\label{eq:tildew-twisted-commutation}
		\tilde{X}^j(\tilde{u})\, \tilde{Z}^{j'}(\tilde{v})
		= (-1)^{\tr((e_j \tilde{u}) \cdot (e_{j'} \tilde{v}))}\, \tilde{Z}^{j'}(\tilde{v})\, \tilde{X}^j(\tilde{u})\;,
	\end{equation}
	where $\tr((e_j \tilde{u}) \cdot (e_{j'} \tilde{v})) = \tr(e_j e_{j'} (\tilde{u} \cdot \tilde{v}))$.
\end{lemma}
\begin{proof}
	\uses{def:tilde-w-observables, def:tilde-m-measurement, def:tau-dot-product-projector, def:s-w-marginals, lem:pauli-observable-expansion, lem:twisted-commutation, lem:qld-4-7}
	Expanding $\tau^W(e_j \tilde{u}) = \sum_{h} (-1)^{\tr(e_j (\tilde{u} \cdot h))}\, \tau^W_h$ on the right-hand side of \eqref{eq:tildewj-product-form} (Lemma~\ref{lem:pauli-observable-expansion}), grouping the $h$ according to the value $a' = h \cdot \tilde{u}$ into the projectors $\tau^W_{a'}(\tilde{u})$ of \eqref{eq:def-tauwu}, collecting the pairs $(g_X,g_Z)$ with a common value $g_W = g$ into $\hat{S}^W_g$ (Definition~\ref{def:s-w-marginals}), and changing the summation variable to $a = (\mathrm{Dec}(g) \cdot \tilde{u}) - a'$ recovers, by \eqref{eq:tilde_M} and \eqref{eq:def-tildewj}, the operator $\tilde{W}^j(\tilde{u})$; this proves \eqref{eq:tildewj-product-form}. Since the $\hat{S}_{g_X,g_Z}$ are mutually orthogonal projectors summing to the identity (Lemma~\ref{lem:qld-4-7}) and $\tau^W(e_j \tilde{u})$ is a Hermitian involution, the product form exhibits $\tilde{W}^j(\tilde{u})$ as Hermitian with square $\sum_{g_X,g_Z} \hat{S}_{g_X,g_Z} \ot \Id = \Id$. For \eqref{eq:tildew-twisted-commutation}, the two sign-weighted sums of the projectors $\hat{S}_{g_X,g_Z}$ appearing in the product forms of $\tilde{X}^j(\tilde{u})$ and $\tilde{Z}^{j'}(\tilde{v})$ commute with each other, so the commutation phase is that of $\tau^X(e_j \tilde{u})$ and $\tau^Z(e_{j'} \tilde{v})$, which is $(-1)^{\tr((e_j \tilde{u}) \cdot (e_{j'} \tilde{v}))}$ by Lemma~\ref{lem:twisted-commutation}.
\end{proof}

% Local correction of the source (lines 1451-1456 of
% 14_analysis_of_the_pauli_basis_test.tex): the source asserts commutation
% whenever j != j'.
\begin{remark}[The commutation phase across basis indices]\label{rem:qld-cross-phase}
	For $j \neq j'$ the source asserts that $\tilde{X}^j(\tilde{u})$ and $\tilde{Z}^{j'}(\tilde{v})$ commute for all $\tilde{u}, \tilde{v} \in \F_q^M$. Since $e_j e_{j'} \neq 0$ and the bilinear form $(a,b) \mapsto \tr(ab)$ on $\F_q$ is nondegenerate, there exist $\tilde{u}, \tilde{v}$ with $\tr(e_j e_{j'} (\tilde{u} \cdot \tilde{v})) = 1$, and for these \eqref{eq:tildew-twisted-commutation} yields anticommutation instead. By self-duality of the basis $\{e_j\}$, the dichotomy asserted by the source --- commutation for $j \neq j'$, and the phase $(-1)^{\tr((e_j\tilde{u}) \cdot (e_j\tilde{v}))}$ for $j = j'$ --- agrees with \eqref{eq:tildew-twisted-commutation} whenever $\tilde{u} \cdot \tilde{v} \in \{0,1\}$, in particular for all $\tilde{u}, \tilde{v} \in \{0,1\}^M$. The relations \eqref{eq:tildew-twisted-commutation} are not used in the remainder of the chapter.
\end{remark}

\section{Consistency with the points measurements}

The next lemma shows that the pulled-apart measurements are consistent with the original points measurements of the strategy, and self-consistent. In its statement, and for the remainder of the chapter, $\delta_S$ denotes the function of Lemma~\ref{lem:qld-4-7}, enlarged by an adjustment of the universal constants $a$ and $b$ so as to also dominate additional error terms of order $\sqrt{\eps}$, $\eps$, and $md/q$ that the source absorbs into $\delta_S$ along the way; such an adjustment is possible since, after shrinking $b$ to at most $1/2$, the term $\eps^b$ dominates $\sqrt{\eps}$ and $\eps$ for $\eps \leq 1$, while $(md)^a q^{-b}$ dominates $md/q$.

\begin{lemma}[Consistency and self-consistency of the pulled-apart measurements]\label{lem:qld-construct-the-paulis}
	\uses{def:tilde-m-measurement, def:tilde-w-observables, lem:projective-strategy-setup, def:expanded-state, def:consistency, def:povm-distance, lem:qld-4-7, def:indicator-vector}
	The measurements $\{\tilde{M}^{W,\tilde{u}}_a\}_{a \in \F_q}$ and observables $\tilde{W}^j(\tilde{u})$ satisfy the following properties.
	% The paper's statement also quantifies item 1 over j, but j does not occur
	% in item 1; the vacuous quantifier is dropped here.
	\begin{enumerate}
		\item (\textbf{Consistency with points measurements}) For all $W \in \{X,Z\}$, on average over uniformly random $u \in \F_q^m$,
		\[
			\Id_{\mathrm{A}\mathrm{A}'\mathrm{A}''} \ot \tilde{M}^{W,\mathrm{ind}_m(u)}_a \simeq_{\delta_S} M^{(\mathrm{Point},W),u}_a \ot \Id_{\mathrm{B}\mathrm{B}'\mathrm{B}''}
		\]
		and
		\[
			\tilde{M}^{W,\mathrm{ind}_m(u)}_a \ot \Id_{\mathrm{B}\mathrm{B}'\mathrm{B}''} \simeq_{\delta_S} \Id_{\mathrm{A}\mathrm{A}'\mathrm{A}''} \ot M^{(\mathrm{Point},W),u}_a\;,
		\]
		with the answer summation over $a \in \F_q$.
		\item (\textbf{Self-consistency}) For all $W \in \{X,Z\}$ and $j \in \{1,\ldots,t\}$, on average over uniformly random $\tilde{u} \in \F_q^{M}$,
		\[
			\tilde{W}^j(\tilde{u}) \ot \Id_{\mathrm{B}\mathrm{B}'\mathrm{B}''} \approx_{\delta_S} \Id_{\mathrm{A}\mathrm{A}'\mathrm{A}''} \ot \tilde{W}^j(\tilde{u})\;.
		\]
	\end{enumerate}
\end{lemma}
\begin{proof}
	\uses{lem:qld-4-7, lem:qld-constructing-the-paulis-helper, def:expanded-point-measurement, def:tilde-m-measurement, def:tau-dot-product-projector, def:decoding-map, def:indicator-vector, lem:qld-win-implications, lem:qld-comm-cons, fact:add-a-proj, fact:agreement, fact:triangle, fact:data-processing, def:bracket, lem:povm-to-obs, lem:schwartz-zippel-total-degree}
	% Support restriction added: the source (line 1492 of
	% 14_analysis_of_the_pauli_basis_test.tex) rewrites Dec(g).ind_m(u) as
	% g(u) for every outcome g; see Remark rem:qld-decoding-identity.
	We first bound the mass that the marginal measurements $\{\hat{S}^W_g\}$ place on polynomials that are not encodings. On average over uniformly random $u \in \F_q^m$, the point-consistency relation \eqref{eq:qld-s-point-con-bob} of Lemma~\ref{lem:qld-4-7} gives $(\hat{S}^W_{[g(u)=a]})_{\mathrm{B}\mathrm{B}'} \approx_{\delta_S} (\hat{M}^{(\mathrm{Point},W),u}_a)_{\mathrm{A}\mathrm{B}''}$ (Lemma~\ref{fact:agreement}); the consistency and Pauli-basis-consistency checks of Lemma~\ref{lem:qld-win-implications} give, again on average over uniform $u$ and combining through Lemma~\ref{fact:data-processing}, Lemma~\ref{fact:agreement}, and Lemma~\ref{fact:triangle}, that $(M^{(\mathrm{Point},W),u}_{a'})_{\mathrm{A}} \approx_{O(\eps)} (M^{(\mathrm{Pauli},W)}_{[g_h(u)=a']})_{\mathrm{A}}$; and replacing the point factor of the convolution form of $\hat{M}^{(\mathrm{Point},W),u}_a$ (Definition~\ref{def:expanded-point-measurement}) accordingly costs the same error, because the projectors $\tau^W_{a''}(\mathrm{ind}_m(u))$ of \eqref{eq:def-tauwu} are mutually orthogonal (Lemma~\ref{fact:add-a-proj}). Writing
	\[
		K^u_a = \sum_{a' + a'' = a} \bigl(M^{(\mathrm{Pauli},W)}_{[g_h(u)=a']}\bigr)_{\mathrm{A}} \ot \bigl(\tau^W_{a''}(\mathrm{ind}_m(u))\bigr)_{\mathrm{B}''}\;,
	\]
	a projective measurement, Lemma~\ref{fact:triangle} and Lemma~\ref{fact:agreement} therefore give, using $\delta_S \geq \eps$,
	\[
		\E_{u \in \F_q^m} \sum_a \bra{\hat\psi} (\hat{S}^W_{[g(u)=a]})_{\mathrm{B}\mathrm{B}'} \ot (K^u_a)_{\mathrm{A}\mathrm{B}''} \ket{\hat\psi} \geq 1 - O(\delta_S)\;.
	\]
	Expanding $K^u_a$ over pairs $(h, h'')$ of Pauli outcomes, a term of the left-hand side with outcomes $(g, h, h'')$ is constrained by $g(u) = g_h(u) + g_{h''}(u) = g_{h + h''}(u)$, by the linearity of the encoding \eqref{eq:low-degree-encoding-definition}. If $g$ is not an encoding then $g \neq g_{h+h''}$, and both are polynomials of total degree at most $md$, so the constraint holds for at most an $md/q$ fraction of $u$ (Lemma~\ref{lem:schwartz-zippel-total-degree}). Since all terms are nonnegative and, for each $g$, the operators of the $(h,h'')$-sum add up to the identity, it follows that
	\begin{equation}\label{eq:qld-nonencoding-mass}
		\sum_{g \,\text{not an encoding}} \bra{\hat\psi} (\hat{S}^W_g)_{\mathrm{B}\mathrm{B}'} \ket{\hat\psi} \,\leq\, O(\delta_S) + \frac{md}{q} \,\leq\, O(\delta_S)\;,
	\end{equation}
	and symmetrically with the registers $\mathrm{A}\mathrm{A}'$ in place of $\mathrm{B}\mathrm{B}'$. (If $d = 0$, every element of $\mathrm{ideg}_{d,m}(\F_q)$ is constant, hence an encoding, and the left-hand side vanishes.)

	For the first item, in the case where $M^{(\mathrm{Point},W),u}_a$ acts on register $\mathrm{A}$ and $\tilde{M}^{W,\mathrm{ind}_m(u)}_a$ on $\mathrm{B}\mathrm{B}'\mathrm{B}''$, expanding \eqref{eq:tilde_M} and \eqref{eq:def-tauwu} and using the convolution form of the combined projectors $\hat{M}^{(\mathrm{Point},W),u}_a$ (Definition~\ref{def:expanded-point-measurement}) gives the exact identity
	\[
		\E_{u \in \F_q^m} \sum_a \bra{\hat\psi} M^{(\mathrm{Point},W),u}_a \ot \tilde{M}^{W,\mathrm{ind}_m(u)}_a \ket{\hat\psi}
		= \E_{u \in \F_q^m} \sum_g \bra{\hat\psi} (\hat{M}^{(\mathrm{Point},W),u}_{\mathrm{Dec}(g) \cdot \mathrm{ind}_m(u)})_{\mathrm{A}\mathrm{B}''} \ot (\hat{S}^W_g)_{\mathrm{B}\mathrm{B}'} \ket{\hat\psi}\;.
	\]
	For encodings $g$ the subscript $\mathrm{Dec}(g) \cdot \mathrm{ind}_m(u)$ equals $g(u)$; discarding the (nonnegative) terms with $g$ not an encoding and reinstating the corresponding terms with subscript $g(u)$, each of which is at most $\bra{\hat\psi} (\hat{S}^W_g)_{\mathrm{B}\mathrm{B}'} \ket{\hat\psi}$, shows that the right-hand side is at least
	\[
		\E_{u \in \F_q^m} \sum_g \bra{\hat\psi} (\hat{M}^{(\mathrm{Point},W),u}_{g(u)})_{\mathrm{A}\mathrm{B}''} \ot (\hat{S}^W_g)_{\mathrm{B}\mathrm{B}'} \ket{\hat\psi} - O(\delta_S)
		\geq 1 - O(\delta_S)
	\]
	by \eqref{eq:qld-nonencoding-mass} and Lemma~\ref{lem:qld-4-7}, and the register-interchanged relation follows analogously. For the second item one first shows that, on average over uniformly random $\tilde{u} \in \F_q^M$,
	\begin{equation}\label{eq:qld-pulling-cons}
		\bigl( \tilde{M}^{W,\tilde{u}}_a \bigr)_{\mathrm{A}\mathrm{A}'\mathrm{A}''} \approx_{\delta_S} \bigl( \tilde{M}^{W,\tilde{u}}_a \bigr)_{\mathrm{B}\mathrm{B}'\mathrm{B}''}\;:
	\end{equation}
	right-multiplying by the resolution $\sum_g (\hat{S}^W_g)_{\mathrm{A}\mathrm{A}'} \ot (\hat{M}^{(\mathrm{Point},W),u}_{g(u)})_{\mathrm{B}\mathrm{A}''} \approx_{\delta_S} \Id$ of Lemma~\ref{lem:qld-constructing-the-paulis-helper} (with Lemma~\ref{fact:add-a-proj}) and expanding the products of $\tau^W$ projectors via \eqref{eq:def-tauwu} and $h \cdot \mathrm{ind}_m(u) = g_h(u)$ places $(\tilde{M}^{W,\tilde{u}}_a)_{\mathrm{A}\mathrm{A}'\mathrm{A}''}$ within $O(\delta_S)$ of a sum over pairs $(g,h)$ with $(\mathrm{Dec}(g)-h)\cdot\tilde{u} = a$. One then applies, in turn, the self-consistency of the points measurements (Lemma~\ref{lem:qld-win-implications}) at cost $\eps$; the exact identity $\ket{\hat\psi} = \sum_{h'} (\tau^W_{h'})_{\mathrm{B}'} \ot (\tau^W_{h'})_{\mathrm{B}''} \ket{\hat\psi}$ of the EPR registers; Lemma~\ref{lem:qld-constructing-the-paulis-helper} twice more (once through a Cauchy--Schwarz detour of cost $O(\sqrt{\eps})$ using the self-consistency of the $\hat{M}^{(\mathrm{Point},W),u}_a$, Lemma~\ref{lem:qld-comm-cons}); and finally a Schwartz--Zippel step (Lemma~\ref{lem:schwartz-zippel-total-degree}, for polynomials of total degree at most $md$) which replaces the event $(g'-g_{h'})(u) = (g-g_h)(u)$ by the polynomial identity $g'-g_{h'} = g-g_h$ at cost $md/q$. The resulting expression
	\[
		\sum_{\substack{g,h,g',h': \\ g' - g_{h'} = g - g_h \\ (\mathrm{Dec}(g)-h)\cdot \tilde{u} = a}}
		\Bigl((\hat{S}^W_g)_{\mathrm{A}\mathrm{A}'} \ot (\tau_h^W)_{\mathrm{A}''} \Bigr) \cdot \Bigl( (\hat{S}^W_{g'})_{\mathrm{B}\mathrm{B}'} \ot (\tau_{h'}^W)_{\mathrm{B}''} \Bigr)
	\]
	is symmetric under exchanging $(g,h)$ with $(g',h')$: since $\mathrm{Dec}$ is linear with $\mathrm{Dec}(g_{h}) = h$ for all $h \in \F_q^M$ (Remark~\ref{rem:qld-decoding-identity}), the constraint $g' - g_{h'} = g - g_h$ implies $\mathrm{Dec}(g) - h = \mathrm{Dec}(g') - h'$, so the two constraint sets coincide. The same derivation therefore bounds the distance of the displayed expression to $(\tilde{M}^{W,\tilde{u}}_a)_{\mathrm{B}\mathrm{B}'\mathrm{B}''}$, and \eqref{eq:qld-pulling-cons} follows after absorbing the error terms $O(\sqrt{\eps})$ and $O(md/q)$ into $\delta_S$. Finally, since $\{\tilde{M}^{W,\tilde{u}}_a\}$ is projective, Item~2 of Lemma~\ref{fact:agreement}, data processing along the map $a \mapsto \tr(e_j a)$ (Lemma~\ref{fact:data-processing}), and Item~1 of Lemma~\ref{fact:agreement} yield $(\tilde{M}^{W,\tilde{u}}_{[\tr(e_j \cdot) = b]})_{\mathrm{A}\mathrm{A}'\mathrm{A}''} \approx_{\delta_S} (\tilde{M}^{W,\tilde{u}}_{[\tr(e_j \cdot) = b]})_{\mathrm{B}\mathrm{B}'\mathrm{B}''}$ with answer summation over $b \in \F_2$, and Lemma~\ref{lem:povm-to-obs} converts this into the claimed self-consistency of the observables $\tilde{W}^j(\tilde{u})$.
\end{proof}

% Following the source, the next lemma is stated after
% lem:qld-construct-the-paulis, whose proof invokes it; its own proof
% depends only on lem:qld-4-7 and lem:qld-comm-cons, so there is no
% circularity.
Lemmas~\ref{lem:qld-4-7} and~\ref{lem:qld-comm-cons} imply the following two estimates, which were used in the proof of Lemma~\ref{lem:qld-construct-the-paulis}.

\begin{lemma}[The marginals absorb the point measurements]\label{lem:qld-constructing-the-paulis-helper}
	\uses{def:s-w-marginals, def:expanded-point-measurement, def:expanded-state, def:povm-distance, lem:qld-4-7, def:symmetric-equivalents}
	For $W \in \{X,Z\}$, on average over uniformly random $u \in \F_q^m$,
	\begin{equation}\label{eq:qld-sg-cons}
		\sum_g (\hat{S}^W_g)_{\mathrm{A}\mathrm{A}'} \ot (\hat{M}^{(\mathrm{Point},W),u}_{g(u)})_{\mathrm{B}\mathrm{A}''} \approx_{\delta_S} \Id
	\end{equation}
	and
	\begin{equation}\label{eq:qld-sg-cons2}
		\bigl(\hat{S}^W_g \cdot (\Id - \hat{M}^{(\mathrm{Point}, W), u}_{g(u)})\bigr)_{\mathrm{A}\mathrm{A}'} \approx_{\delta_S} 0\;,
	\end{equation}
	where the answer summation is over $g \in \mathrm{ideg}_{d,m}(\F_q)$, and where in \eqref{eq:qld-sg-cons2} an additional error term of order $\eps$ is absorbed into $\delta_S$. Furthermore, all symmetric equivalents of these approximations also hold (Definition~\ref{def:symmetric-equivalents}).
\end{lemma}
\begin{proof}
	\uses{lem:qld-4-7, fact:agreement, fact:add-a-proj, def:bracket, lem:qld-comm-cons}
	The point-consistency conclusion of Lemma~\ref{lem:qld-4-7} states that, on average over $u \in \F_q^m$,
	\[
		\bra{\hat\psi} \sum_g (\hat{S}^W_g)_{\mathrm{A}\mathrm{A}'} \ot (\hat{M}^{(\mathrm{Point}, W), u}_{g(u)})_{\mathrm{B}\mathrm{A}''} \ket{\hat\psi} \geq 1 - \delta_S\;,
	\]
	which says exactly that the operator on the left-hand side of \eqref{eq:qld-sg-cons} is $\delta_S$-consistent with $\Id$; Lemma~\ref{fact:agreement} upgrades this $\simeq_{\delta_S}$ statement to the $\approx_{\delta_S}$ statement \eqref{eq:qld-sg-cons}. For \eqref{eq:qld-sg-cons2}, the same consistency, written with the coarse-graining $\hat{S}^W_{[g(u)=a]} = \sum_{g : g(u) = a} \hat{S}^W_g$ as $(\hat{S}^W_{[g(u)=a]})_{\mathrm{A}\mathrm{A}'} \approx_{\delta_S} (\hat{M}^{(\mathrm{Point}, W), u}_{a})_{\mathrm{B}\mathrm{A}''}$ with answer summation over $a \in \F_q$, is left-multiplied by $(\hat{S}^W_{[g(u)=a]})_{\mathrm{A}\mathrm{A}'}$ using Lemma~\ref{fact:add-a-proj}; moving $\hat{M}^{(\mathrm{Point}, W), u}_a$ from registers $\mathrm{B}\mathrm{A}''$ to $\mathrm{A}\mathrm{A}'$ costs a further $\eps$ by Cauchy--Schwarz together with the self-consistency of the $\hat{M}^{(\mathrm{Point}, W), u}_a$ (Lemma~\ref{lem:qld-comm-cons}), giving
	\[
		(\hat{S}^W_{[g(u)=a]})_{\mathrm{A}\mathrm{A}'} \approx_{O(\delta_S + \eps)} (\hat{S}^W_{[g(u)=a]} \cdot \hat{M}^{(\mathrm{Point}, W), u}_{a})_{\mathrm{A}\mathrm{A}'}\;.
	\]
	Expanding
	\[
		\E_u \sum_g \bigl\| (\hat{S}^W_{g} \cdot (\Id - \hat{M}^{(\mathrm{Point}, W), u}_{g(u)}))_{\mathrm{A}\mathrm{A}'} \ket{\hat\psi} \bigr\|^2
		= \E_u \sum_a \bra{\hat\psi} \bigl((\Id - \hat{M}^{(\mathrm{Point}, W), u}_{a}) \cdot \hat{S}^W_{[g(u)=a]} \cdot (\Id - \hat{M}^{(\mathrm{Point}, W), u}_{a})\bigr)_{\mathrm{A}\mathrm{A}'} \ket{\hat\psi}
	\]
	and bounding the right-hand side by $O(\delta_S + \eps)$ using the two approximations above yields \eqref{eq:qld-sg-cons2} after absorbing $\eps$ into $\delta_S$. The symmetric equivalents follow by the same argument with the register bipartitions interchanged.
\end{proof}

\section{The swap unitaries and the extraction of the Pauli group}

\begin{definition}[The swap unitaries]\label{def:v-swap-unitary}
	\uses{lem:qld-4-7, def:generalized-pauli, def:decoding-map}
	Define the linear map
	\[
		V_{\mathrm{A}} = \sum_{g_X,g_Z} (\hat{S}_{g_X,g_Z})_{\mathrm{A}\mathrm{A}'} \ot \bigl(\tau^X(\mathrm{Dec}(g_Z)) \, \tau^Z(\mathrm{Dec}(g_X))\bigr)_{\mathrm{A}''}\;,
	\]
	where $\tau^X(\cdot)$ and $\tau^Z(\cdot)$ are the generalized Pauli observables acting on $(\C^q)^{\ot M}$; note that the argument of $\tau^X$ is $\mathrm{Dec}(g_Z)$ and the argument of $\tau^Z$ is $\mathrm{Dec}(g_X)$. The map $V_{\mathrm{B}}$, acting on registers $\mathrm{B}\mathrm{B}'\mathrm{B}''$, is defined analogously.
\end{definition}

% In the source the following properties are established inline with the
% definition of the swap unitaries; they are recorded here as a lemma.
\begin{lemma}[Unitarity and conjugation action of the swap unitaries]\label{lem:v-swap-conjugation}
	\uses{def:v-swap-unitary, def:tilde-w-observables, def:tilde-m-measurement, def:generalized-pauli, def:indicator-vector, def:bracket}
	The maps $V_{\mathrm{A}}$ and $V_{\mathrm{B}}$ of Definition~\ref{def:v-swap-unitary} are unitaries acting on registers $\mathrm{A}\mathrm{A}'\mathrm{A}''$ and $\mathrm{B}\mathrm{B}'\mathrm{B}''$ respectively, and their conjugation action is exact, as follows. For all $W \in \{X,Z\}$, $j \in \{1,\ldots,t\}$, and $\tilde{u} \in \F_q^M$,
	\begin{equation}\label{eq:v-swap-obs-conjugation}
		V_{\mathrm{A}} \, \tilde{W}^j(\tilde{u}) \, V_{\mathrm{A}}^\dagger = \Id_{\mathrm{A}\mathrm{A}'} \ot \bigl( \tau^W(e_j \tilde{u}) \bigr)_{\mathrm{A}''}
		\qquad \text{and} \qquad
		V_{\mathrm{B}} \, \tilde{W}^j(\tilde{u}) \, V_{\mathrm{B}}^\dagger = \Id_{\mathrm{B}\mathrm{B}'} \ot \bigl( \tau^W(e_j \tilde{u}) \bigr)_{\mathrm{B}''}\;,
	\end{equation}
	and for all $W \in \{X,Z\}$, $u \in \F_q^m$, and $a \in \F_q$,
	\begin{equation}\label{eq:qld-unitary-6}
		V_{\mathrm{B}} \, \tilde{M}^{W,\mathrm{ind}_m(u)}_a \, V_{\mathrm{B}}^\dagger = \Id_{\mathrm{B}\mathrm{B}'} \ot (\tau^W_{[g_h(u)=a]})_{\mathrm{B}''}\;,
		\qquad
		\tau^W_{[g_h(u)=a]} = \sum_{h \,:\, g_h(u) = a} \tau^W_h\;,
	\end{equation}
	together with the analogous identity for $V_{\mathrm{A}}$ on registers $\mathrm{A}\mathrm{A}'\mathrm{A}''$.
\end{lemma}
\begin{proof}
	\uses{lem:qld-4-7, lem:tildew-product-form, lem:twisted-commutation, def:v-swap-unitary, def:tilde-m-measurement, def:tau-dot-product-projector, def:indicator-vector, def:decoding-map, def:generalized-pauli}
	Since the measurement $\{\hat{S}_{g_X,g_Z}\}$ is projective (Lemma~\ref{lem:qld-4-7}) and the $\tau^X, \tau^Z$ observables are self-inverse,
	\[
		V_{\mathrm{A}} V_{\mathrm{A}}^\dagger = \sum_{g_X,g_Z} (\hat{S}_{g_X,g_Z})_{\mathrm{A}\mathrm{A}'} \ot \Id_{\mathrm{A}''} = \Id_{\mathrm{A}\mathrm{A}'\mathrm{A}''}\;,
	\]
	and the same computation gives $V_{\mathrm{A}}^\dagger V_{\mathrm{A}} = \Id$, so $V_{\mathrm{A}}$ is unitary; likewise $V_{\mathrm{B}}$. For \eqref{eq:v-swap-obs-conjugation}, conjugate the product form \eqref{eq:tildewj-product-form} of $\tilde{W}^j(\tilde{u})$ (Lemma~\ref{lem:tildew-product-form}): by the orthogonality of the projectors $\hat{S}_{g_X,g_Z}$, only the diagonal terms survive, and the twisted commutation relation (Lemma~\ref{lem:twisted-commutation})
	\[
		\tau^X(\mathrm{Dec}(g_Z)) \, \tau^Z(\mathrm{Dec}(g_X)) \, \tau^W(e_j \tilde{u}) = (-1)^{\tr( e_j(\mathrm{Dec}(g_W) \cdot \tilde{u}))} \, \tau^W(e_j \tilde{u}) \, \tau^X(\mathrm{Dec}(g_Z)) \, \tau^Z(\mathrm{Dec}(g_X))
	\]
	cancels the sign $(-1)^{\tr(e_j (\mathrm{Dec}(g_W) \cdot \tilde{u}))}$ term by term, leaving $\sum_{g_X,g_Z} \hat{S}_{g_X,g_Z} \ot \tau^W(e_j \tilde{u}) = \Id \ot \tau^W(e_j \tilde{u})$. For \eqref{eq:qld-unitary-6}, expand $\tilde{M}^{W,\mathrm{ind}_m(u)}_a$ by \eqref{eq:tilde_M} and \eqref{eq:def-tauwu}; by the orthogonality of the $\hat{S}_{g_X,g_Z}$, conjugation replaces, in the term of $\hat{S}^W_g$, each projector $\tau^W_h$ by $\tau^W_{h + \mathrm{Dec}(g)}$, using the identities $\tau^X(s)\, \tau^Z_h \,\tau^X(s)^\dagger = \tau^Z_{h+s}$ and $\tau^Z(s)\, \tau^X_h\, \tau^Z(s)^\dagger = \tau^X_{h+s}$. Substituting $h' = h + \mathrm{Dec}(g)$ turns the constraint $h \cdot \mathrm{ind}_m(u) = (\mathrm{Dec}(g) \cdot \mathrm{ind}_m(u)) - a$ into $h' \cdot \mathrm{ind}_m(u) = -a = a$, since $\F_q$ has characteristic $2$; the quantity $\mathrm{Dec}(g) \cdot \mathrm{ind}_m(u)$ cancels exactly, so no relation between it and $g(u)$ is used (Remark~\ref{rem:qld-decoding-identity}). Since $h' \cdot \mathrm{ind}_m(u) = g_{h'}(u)$ \eqref{eq:low-degree-encoding-definition}, summing over $g$ gives \eqref{eq:qld-unitary-6}; the computation for $V_{\mathrm{A}}$ on registers $\mathrm{A}\mathrm{A}'\mathrm{A}''$ is identical.
\end{proof}

\begin{lemma}[Extraction of the EPR state and the Pauli observables]\label{lem:qld-unitary}
	\uses{lem:qld-4-7, def:v-swap-unitary, def:expanded-state, lem:projective-strategy-setup, def:generalized-pauli, def:EPR, def:povm-distance}
	There exists a function $\delta_{\mathrm{qld}}(\eps,m,d,q) = O(\delta_S^{1/4} + md/q)$, where $\delta_S$ is the function of Lemma~\ref{lem:qld-4-7}, and there exist unitaries $V_{\mathrm{A}}$ acting on registers $\mathrm{A}\mathrm{A}'\mathrm{A}''$ and $V_{\mathrm{B}}$ acting on registers $\mathrm{B}\mathrm{B}'\mathrm{B}''$, such that the following hold.
	\begin{enumerate}
		\item There exists a state $\ket{\mathrm{aux}}$ on registers $\mathrm{A}\mathrm{A}'\mathrm{B}\mathrm{B}'$ such that
		\[
			\bigl\| V_{\mathrm{A}} \ot V_{\mathrm{B}} \ket{\hat{\psi}}_{\mathrm{A}\mathrm{A}'\mathrm{A}''\mathrm{B}\mathrm{B}'\mathrm{B}''} - \ket{\mathrm{aux}}_{\mathrm{A}\mathrm{A}'\mathrm{B}\mathrm{B}'} \ot \ket{\mathrm{EPR}_q}^{\ot M}_{\mathrm{A}''\mathrm{B}''} \bigr\|^2 \leq \delta_{\mathrm{qld}}\;.
		\]
		\item For $W \in \{X,Z\}$,
		\begin{align*}
			V_{\mathrm{A}} \, \bigl(M^{(\mathrm{Pauli},W)}_h \bigr)_{\mathrm{A}} \, V_{\mathrm{A}}^\dagger &\approx_{\delta_{\mathrm{qld}}} \Id_{\mathrm{A}\mathrm{A}'} \ot (\tau^W_h)_{\mathrm{A}''}\;, \\
			V_{\mathrm{B}} \, \bigl(M^{(\mathrm{Pauli},W)}_h \bigr)_{\mathrm{B}} \, V_{\mathrm{B}}^\dagger &\approx_{\delta_{\mathrm{qld}}} \Id_{\mathrm{B}\mathrm{B}'} \ot (\tau^W_h)_{\mathrm{B}''}\;,
		\end{align*}
		where the $\approx$ statements hold with respect to the state $\ket{\mathrm{aux}}_{\mathrm{A}\mathrm{A}'\mathrm{B}\mathrm{B}'} \ot \ket{\mathrm{EPR}_q}^{\ot M}_{\mathrm{A}''\mathrm{B}''}$, with the answer summation over $h \in \F_q^{M}$ and no average over questions.
	\end{enumerate}
	The unitaries may be taken to be $V_{\mathrm{A}}$ and $V_{\mathrm{B}}$ of Definition~\ref{def:v-swap-unitary}.
\end{lemma}
\begin{proof}
	\uses{def:v-swap-unitary, lem:v-swap-conjugation, lem:tildew-product-form, lem:qld-construct-the-paulis, lem:qld-win-implications, fact:triangle-for-simeq, def:bracket, lem:cancellation, lem:schwartz-zippel-total-degree, lem:qld-4-7}
	For the first item, the self-consistency of the observables $\tilde{W}^j$ (Item~2 of Lemma~\ref{lem:qld-construct-the-paulis}) is equivalent, since these observables are Hermitian involutions (Lemma~\ref{lem:tildew-product-form}), to $\E_{j \in \{1,\ldots,t\}} \E_{\tilde{u} \in \F_q^M} \bra{\hat\psi} \tilde{W}^j(\tilde{u}) \ot \tilde{W}^j(\tilde{u}) \ket{\hat\psi} \geq 1 - \delta_S/2$; setting $\ket{\vartheta} = V_{\mathrm{A}} \ot V_{\mathrm{B}} \ket{\hat\psi}$ and using the exact conjugation identities of Lemma~\ref{lem:v-swap-conjugation}, this becomes $\bra{\vartheta} H_W \ket{\vartheta} \geq 1 - \delta_S/2$ for $W \in \{X,Z\}$, where $H_W = \E_j \E_{\tilde{u}} \, \tau^W(e_j \tilde{u}) \ot \tau^W(e_j \tilde{u}) = \bigl( \E_{s \in \F_q} \tau^W(s) \ot \tau^W(s) \bigr)^{\ot M}$ by the tensor factorization of $\tau^W(\tilde{u})$ and the invariance of the uniform average over $s$ under $s \mapsto e_j s$. By Lemma~\ref{lem:cancellation}, $\E_{s,s'} \tau^X(s)\tau^Z(s') \ot \tau^X(s)\tau^Z(s') = \ket{\mathrm{EPR}_q}\bra{\mathrm{EPR}_q}$ on each pair of qudits, whence $H_X H_Z = \ket{\mathrm{EPR}_q}\bra{\mathrm{EPR}_q}^{\ot M}$. The bounds $\| (\Id - H_W)\ket{\vartheta} \|^2 \leq \delta_S$ and $\bra{\vartheta} H_W^2 \ket{\vartheta} \geq 1 - 2\delta_S$, combined through the triangle inequality $\| H_X \ket{\vartheta} - H_Z \ket{\vartheta} \| \leq 2\sqrt{\delta_S}$ (see Remark~\ref{rem:qld-unitary-triangle-slip}), give $\bra{\vartheta} \Id_{\mathrm{A}\mathrm{A}'\mathrm{B}\mathrm{B}'} \ot (\ket{\mathrm{EPR}_q}\bra{\mathrm{EPR}_q}^{\ot M})_{\mathrm{A}''\mathrm{B}''} \ket{\vartheta} \geq 1 - 4\delta_S$.
	% Case distinction added: the source normalizes the projected vector
	% without ensuring that it is nonzero; see Remark rem:qld-unitary-triangle-slip.
	If $\delta_S > 1/8$, then, after enlarging the universal constant in $\delta_{\mathrm{qld}}$, both items of the lemma hold with $\ket{\mathrm{aux}}$ an arbitrary unit state, since their left-hand sides never exceed $4$; assume therefore $\delta_S \leq 1/8$. Then the vector $(\bra{\mathrm{EPR}_q}^{\ot M})_{\mathrm{A}''\mathrm{B}''} \ket{\vartheta}$ has squared norm at least $1 - 4\delta_S \geq 1/2 > 0$, and $\ket{\mathrm{aux}}$ may be defined as its normalization; the overlap of $\ket{\vartheta}$ with $\ket{\mathrm{aux}} \ot \ket{\mathrm{EPR}_q}^{\ot M}$ equals the norm of that vector and is at least $\sqrt{1 - 4\delta_S} \geq 1 - 4\delta_S$, which yields the first item with a bound of $8\delta_S = O(\sqrt{\delta_S})$. For the second item, Lemma~\ref{lem:qld-win-implications} gives, on average over $u \in \F_q^m$, $M^{(\mathrm{Pauli},W)}_{[g_h(u)=a]} \ot \Id \simeq_\eps \Id \ot M^{(\mathrm{Point},W),u}_a$ and $\Id \ot M^{(\mathrm{Point},W),u}_a \simeq_\eps M^{(\mathrm{Point},W),u}_a \ot \Id$, where $M^{(\mathrm{Pauli},W)}_{[g_h(u)=a]} = \sum_{h : g_h(u) = a} M^{(\mathrm{Pauli},W)}_h$; combining these with Item~1 of Lemma~\ref{lem:qld-construct-the-paulis} by Lemma~\ref{fact:triangle-for-simeq}, using $\delta_S \geq \eps$, yields $M^{(\mathrm{Pauli},W)}_{[g_h(u)=a]} \ot \Id \simeq_{\sqrt{\delta_S}} \Id \ot \tilde{M}^{W,\mathrm{ind}_m(u)}_a$. Writing $\ket{\Delta} = \ket{\mathrm{aux}} \ot \ket{\mathrm{EPR}_q}^{\ot M}$, the quantity to be bounded for the $\mathrm{A}$-side relation equals $2 - 2\sum_h \bra{\Delta} V_{\mathrm{A}} M^{(\mathrm{Pauli},W)}_h V_{\mathrm{A}}^\dagger \ot (\tau^W_h)_{\mathrm{B}''} \ket{\Delta}$, since $(\tau^W_h)_{\mathrm{A}''} \ket{\mathrm{EPR}_q}^{\ot M} = (\tau^W_h)_{\mathrm{B}''} \ket{\mathrm{EPR}_q}^{\ot M}$. Replacing $(\tau^W_h)_{\mathrm{B}''}$ by the averaged coarse-graining $\E_u \sum_{h' : g_{h'}(u) = g_h(u)} (\tau^W_{h'})_{\mathrm{B}''}$ costs at most $md/q$ by Lemma~\ref{lem:schwartz-zippel-total-degree} (distinct $h \neq h'$ have distinct encodings $g_h \neq g_{h'}$, of total degree at most $md$), and by \eqref{eq:qld-unitary-6} the resulting sum equals $\E_u \sum_{a} \bra{\Delta} V_{\mathrm{A}} M^{(\mathrm{Pauli},W)}_{[g_h(u)=a]} V_{\mathrm{A}}^\dagger \ot V_{\mathrm{B}} \tilde{M}^{W,\mathrm{ind}_m(u)}_a V_{\mathrm{B}}^\dagger \ket{\Delta}$. Exchanging $\ket{\Delta}$ for $\ket{\vartheta}$ costs $O(\delta_S^{1/4})$ by the first item, and the resulting expression $\E_u \sum_a \bra{\hat\psi} M^{(\mathrm{Pauli},W)}_{[g_h(u)=a]} \ot \tilde{M}^{W,\mathrm{ind}_m(u)}_a \ket{\hat\psi}$ is at least $1 - O(\sqrt{\delta_S})$ by the $\simeq_{\sqrt{\delta_S}}$ relation above. Altogether the distance is $O(\delta_S^{1/4} + md/q)$; the $\mathrm{B}$-side relation is proved identically.
\end{proof}

\begin{remark}[Two steps in the source proof of the extraction lemma]\label{rem:qld-unitary-triangle-slip}
	In the source proof of Lemma~\ref{lem:qld-unitary}, the triangle inequality gives $\| H_X \ket{\vartheta} - H_Z \ket{\vartheta} \| \leq \| (\Id - H_X)\ket{\vartheta} \| + \| (\Id - H_Z)\ket{\vartheta} \| \leq 2\sqrt{\delta_S}$, so that $\| H_X \ket{\vartheta} - H_Z \ket{\vartheta} \|^2 \leq 4\delta_S$; the source instead asserts the bound $2\sqrt{\delta_S}$ for the \emph{squared} norm, which holds only when $\sqrt{\delta_S} \leq 1/2$, and consequently carries the term $\sqrt{\delta_S}$ into an overlap bound of $1 - 2\delta_S - \sqrt{\delta_S}$. The proof above uses the corrected bound $4\delta_S$, which gives the overlap bound $1 - 4\delta_S$ and leaves the final bound $\delta_{\mathrm{qld}} = O(\delta_S^{1/4} + md/q)$ unchanged. The source proof also normalizes the vector $(\bra{\mathrm{EPR}_q}^{\ot M})_{\mathrm{A}''\mathrm{B}''} \ket{\vartheta}$ without ensuring that it is nonzero; the case distinction on $\delta_S$ in the proof above supplies this, the conclusions of the lemma being trivial when $\delta_S$ exceeds a universal constant.
\end{remark}

\section{Proof of the soundness of the Pauli basis test}

We can now prove the soundness theorem for the Pauli basis test, stated as Theorem~\ref{thm:pauli} in Chapter~\ref{chap:qpbt-test}.

\begin{proof}
	\proves{thm:pauli}
	\uses{lem:qld-unitary, lem:pauli-completeness, thm:naimark, lem:projective-strategy-setup, def:expanded-state, def:EPR, def:povm-distance}
	Let $(\ket{\psi},M)$ be a strategy for $\mathfrak{G}^{\mathrm{Pauli}}_{(q,m,d)}$ that succeeds with probability at least $1 - \eps$, where $\ket{\psi}$ is a bipartite state on registers $\mathrm{A}$ and $\mathrm{B}$. By Naimark's theorem (Theorem~\ref{thm:naimark}, as formulated in~\cite[Theorem 5.1]{Ji2020LowIndividualDegree}), at the cost of extending $\ket{\psi}$ by sufficiently many ancilla qubits initialized in the state $\ket{0}$, we may assume that the strategy is projective and that the state is padded as in the conventions of Lemma~\ref{lem:projective-strategy-setup} and the convention following it. Let $V_{\mathrm{A}}$ be the unitary acting on registers $\mathrm{A}\mathrm{A}'\mathrm{A}''$ given by Lemma~\ref{lem:qld-unitary}, and let $\phi_{\mathrm{A}}$ be the isometry mapping register $\mathrm{A}$ to registers $\mathrm{A}\mathrm{A}'\mathrm{A}''$, where $\mathrm{A}'$ and $\mathrm{A}''$ are isomorphic to $(\C^q)^{\ot M}$, defined by
	\[
		\phi_{\mathrm{A}}\colon \ket{\theta}_{\mathrm{A}} \mapsto V_{\mathrm{A}} \bigl(\ket{\theta}_{\mathrm{A}} \ot (\ket{\mathrm{EPR}_q}^{\ot M})_{\mathrm{A}' \mathrm{A}''}\bigr)\;;
	\]
	define $\phi_{\mathrm{B}}$ analogously.
	% Reliance disclosure added: the proof chain of thm:pauli passes through
	% lem:qld-4-7, whose source proof (lines 1277--1413 of
	% 14_analysis_of_the_pauli_basis_test.tex) instantiates the classical low
	% individual degree test at dimension 2m+2 without a divisibility
	% hypothesis; the sentence below records the resulting outstanding
	% obligation.
	This invocation is where the analysis of the classical low individual degree test enters: Lemma~\ref{lem:qld-unitary} rests on the polynomial-pair measurement of Lemma~\ref{lem:qld-4-7}, whose proof instantiates the classical test in $2m+2$ variables and at present does so only when $2m+2$ divides $q$ --- a divisibility that fails for every admissible tuple with $m \geq 2$. Three further steps of that same chain are likewise unproved. The classical soundness statement invoked there, Theorem~\ref{lem:ld-soundness}, is imported from Theorem~4.7 of~\cite{Ji2021Quantum} through an identification of the game $\mathfrak{G}^{\mathrm{ld}}_{(q,m,d,1)}$ with the two-prover tensor-code test for $\code^{\ot m}$, and neither that game correspondence nor the bound on the auxiliary parameter of Theorem~4.7 under which the stated error function $\delta_{\mathrm{ld}}$ is obtained has been established; both, with the divisibility obstruction, are documented in \cite{gap:qpbt_ld-dimension-divisibility}. Furthermore, the combined-lines consistency of Lemma~\ref{lem:qld-4-13}, which Lemma~\ref{lem:qld-4-7} consumes, is established here only with the error $m \cdot \poly(\eps, md/q)$, the error form $\poly(m^2 \eps, md/q)$ carried in its statement remaining unproved (see \cite{gap:qpbt_combined-lines-error-term}). The present proof is therefore complete only modulo these four obligations. Since the expanded state $\ket{\hat\psi}$ of Definition~\ref{def:expanded-state} is exactly $\ket{\psi}$ with the halves of $M$ pairs $\ket{\mathrm{EPR}_q}$ appended locally in registers $\mathrm{A}'\mathrm{A}''$ and $\mathrm{B}'\mathrm{B}''$, we have $\phi_{\mathrm{A}} \ot \phi_{\mathrm{B}} \ket{\psi} = V_{\mathrm{A}} \ot V_{\mathrm{B}} \ket{\hat\psi}$. Lemma~\ref{lem:qld-unitary} therefore yields a state $\ket{\mathrm{aux}}$ on registers $\mathrm{A}\mathrm{A}'\mathrm{B}\mathrm{B}'$ such that
	\[
		\bigl\| \phi_{\mathrm{A}} \ot \phi_{\mathrm{B}} \ket{\psi} - \ket{\mathrm{aux}} \ot \ket{\mathrm{EPR}_q}^{\ot M} \bigr\|^2 \leq \delta_{\mathrm{qld}}\;.
	\]
	% Transfer step added: the source (lines 1864--1875 of
	% 14_analysis_of_the_pauli_basis_test.tex) passes directly from the
	% conjugation estimates of lem:qld-unitary, which concern the
	% unitarily conjugated operators V (M_h ot Id) V^dagger, to the
	% corresponding estimates for the isometry phi, whose conjugation
	% inserts the projection onto its range; the range-projection
	% argument below supplies the missing step.
	Item~2 of Lemma~\ref{lem:qld-unitary} concerns the unitarily conjugated operators $N^W_h = V_{\mathrm{A}} \, (M^{(\mathrm{Pauli},W)}_h)_{\mathrm{A}} \, V_{\mathrm{A}}^\dagger$, whereas conjugation by the isometry $\phi_{\mathrm{A}}$ additionally inserts the projection $Q_{\mathrm{A}} = \phi_{\mathrm{A}} \phi_{\mathrm{A}}^\dagger = V_{\mathrm{A}} \, \bigl(\ket{\mathrm{EPR}_q}\bra{\mathrm{EPR}_q}^{\ot M}\bigr)_{\mathrm{A}'\mathrm{A}''} \, V_{\mathrm{A}}^\dagger$ onto the range of $\phi_{\mathrm{A}}$: since $(M^{(\mathrm{Pauli},W)}_h)_{\mathrm{A}}$ and $\bigl(\ket{\mathrm{EPR}_q}\bra{\mathrm{EPR}_q}^{\ot M}\bigr)_{\mathrm{A}'\mathrm{A}''}$ act on disjoint registers and hence commute,
	\[
		\phi_{\mathrm{A}} \, (M^{(\mathrm{Pauli},W)}_h)_{\mathrm{A}} \, \phi_{\mathrm{A}}^\dagger
		= V_{\mathrm{A}} \, \Bigl( (M^{(\mathrm{Pauli},W)}_h)_{\mathrm{A}} \ot \bigl(\ket{\mathrm{EPR}_q}\bra{\mathrm{EPR}_q}^{\ot M}\bigr)_{\mathrm{A}'\mathrm{A}''} \Bigr) \, V_{\mathrm{A}}^\dagger
		= N^W_h \, Q_{\mathrm{A}} = Q_{\mathrm{A}} \, N^W_h\;.
	\]
	The estimates transfer from $N^W_h$ to $\phi_{\mathrm{A}} \, (M^{(\mathrm{Pauli},W)}_h)_{\mathrm{A}} \, \phi_{\mathrm{A}}^\dagger$ as follows. Write $\ket{\Delta} = \ket{\mathrm{aux}} \ot \ket{\mathrm{EPR}_q}^{\ot M}$ and $\ket{\vartheta} = \phi_{\mathrm{A}} \ot \phi_{\mathrm{B}} \ket{\psi}$. Since $\phi_{\mathrm{A}}^\dagger \phi_{\mathrm{A}} = \Id$, the projection $\Id - Q_{\mathrm{A}}$ annihilates $\ket{\vartheta}$, so the displayed bound above gives
	\[
		\bigl\| (\Id - Q_{\mathrm{A}}) \ket{\Delta} \bigr\|^2
		= \bigl\| (\Id - Q_{\mathrm{A}}) (\ket{\Delta} - \ket{\vartheta}) \bigr\|^2
		\leq \bigl\| \ket{\Delta} - \ket{\vartheta} \bigr\|^2
		\leq \delta_{\mathrm{qld}}\;.
	\]
	Since the strategy is projective, the operators $N^W_h$, for $h \in \F_q^M$, are mutually orthogonal projectors summing to the identity, so $\sum_h \| N^W_h \ket{\xi} \|^2 = \| \ket{\xi} \|^2$ for every vector $\ket{\xi}$. Decomposing
	\[
		\bigl( \phi_{\mathrm{A}} \, (M^{(\mathrm{Pauli},W)}_h)_{\mathrm{A}} \, \phi_{\mathrm{A}}^\dagger - (\tau^W_h)_{\mathrm{A}''} \bigr) \ket{\Delta}
		= \bigl( N^W_h - (\tau^W_h)_{\mathrm{A}''} \bigr) \ket{\Delta} - N^W_h (\Id - Q_{\mathrm{A}}) \ket{\Delta}
	\]
	and summing the squared norms over $h$ via $\| \xi - \xi' \|^2 \leq 2\|\xi\|^2 + 2\|\xi'\|^2$, the first terms contribute $O(\delta_{\mathrm{qld}})$ by Item~2 of Lemma~\ref{lem:qld-unitary} and the second terms contribute at most $2 \| (\Id - Q_{\mathrm{A}}) \ket{\Delta} \|^2 \leq 2\delta_{\mathrm{qld}}$; hence, for $W \in \{X,Z\}$,
	\[
		\phi_{\mathrm{A}} \, (M^{(\mathrm{Pauli},W)}_h)_{\mathrm{A}} \, \phi_{\mathrm{A}}^\dagger \approx_{\delta_{\mathrm{qld}}} (\tau^W_h)_{\mathrm{A}''}
		\qquad \text{and, by the identical argument,} \qquad
		\phi_{\mathrm{B}} \, (M^{(\mathrm{Pauli},W)}_h)_{\mathrm{B}} \, \phi_{\mathrm{B}}^\dagger \approx_{\delta_{\mathrm{qld}}} (\tau^W_h)_{\mathrm{B}''}\;,
	\]
	where the $\approx$ statements hold with respect to the state $\ket{\mathrm{aux}} \ot \ket{\mathrm{EPR}_q}^{\ot M}$ and the answer summation is over $h \in \F_q^M$; the constant implicit in $\approx_{\delta_{\mathrm{qld}}}$ (Definition~\ref{def:povm-distance}) absorbs the factors accumulated above. This shows that the game $\mathfrak{G}^{\mathrm{Pauli}}_{(q,m,d)}$ is a self-test, with robustness $\delta_{\mathrm{qld}}(\eps,m,d,q)$, for the canonical Pauli strategy of Lemma~\ref{lem:pauli-completeness}. Identifying the registers $\mathrm{A}\mathrm{A}'$ with $\mathrm{A}'$ and $\mathrm{B}\mathrm{B}'$ with $\mathrm{B}'$ matches the statement of Theorem~\ref{thm:pauli}, in which $\phi_{\mathrm{A}}\colon \mathcal{H}_{\mathrm{A}} \to \mathcal{H}_{\mathrm{A}'} \ot \mathcal{H}_{\mathrm{A}''}$. Finally, unpacking the dependence on the parameters $\eps, m, d, q$, the function $\delta_{\mathrm{qld}}(\eps,m,d,q) = O(\delta_S^{1/4} + md/q)$ has the form $a(md)^a(\eps^b + q^{-b} + 2^{-bmd})$ for universal constants $a > 1$ and $0 < b < 1$; see Remark~\ref{rem:pauli-robustness-form} for the reconciliation of this form with the statement of Theorem~\ref{thm:pauli}.
\end{proof}

\begin{remark}[Norm conventions in the robustness bound]\label{rem:pauli-robustness-form}
	The proof above, following the source, establishes the first conclusion of Theorem~\ref{thm:pauli} in the squared form $\| \phi_{\mathrm{A}} \ot \phi_{\mathrm{B}} \ket{\psi} - \ket{\mathrm{aux}} \ot \ket{\mathrm{EPR}_q}^{\ot M} \|^2 \leq \delta_{\mathrm{qld}}$, whereas the theorem bounds the unsquared norm by $\delta_{\mathrm{qld}}(\eps,m,d,q)$. The two are reconciled by replacing the exponent $b$ with $b/2$ and enlarging $a$, since the class of functions $a(md)^a(\eps^b + q^{-b} + 2^{-bmd})$ is closed under square roots up to such adjustments of the universal constants. Similarly, the closing sentence of the source proof asserts constants $a > 1$ and $0 < b < 1$, while Theorem~\ref{thm:pauli} requires only $a \geq 1$ and $0 < b < 1$; the former implies the latter.
\end{remark}
